\chapter{Unit 6: Innovation Project, Re-Design and Product Presentation}
\label{chap:unit6_qb}

\section*{Practice Question Set (30 MCQs)}
\addcontentsline{toc}{section}{Practice Question Set (30 MCQs)}

\mcqitem{1}{What is "Design Iteration"?}{Copying an existing design without modification}{The cyclical process of prototyping, testing, analyzing feedback, and refining a solution}{Writing code once and never modifying it}{Designing only on alternate days of the week}

\mcqitem{2}{What is the fundamental difference between Design Refactoring and Product Re-Design?}{Refactoring changes the logo color; Re-design deletes the project}{Refactoring improves internal structure and consistency without changing core function; Re-design fundamentally alters structure, form, or workflow}{Refactoring is done by designers; Re-design is done by lawyers}{There is no difference between refactoring and re-design}

\mcqitem{3}{When should a team undertake a fundamental Product Re-Design rather than minor iteration?}{When a single user dislikes a font}{When testing reveals that core problem assumptions, basic architecture, or safety mechanisms have failed}{When the sprint deadline is 2 hours away}{When the marketing team wants a new slogan}

\mcqitem{4}{In Lean Startup methodology, what is a "Strategic Pivot"?}{Rotating a 3D model in CAD software}{A structured course correction designed to test a new fundamental hypothesis about the product, customer segment, or business model while keeping validated learnings}{Firing the entire executive board}{Shutting down the company permanently}

\mcqitem{5}{Which of the following is a classic example of a "Customer Segment Pivot"?}{Rewriting a Python backend in Rust}{Shifting a product originally built for individual college students to enterprise corporate HR departments}{Changing the company office address}{Adding dark mode to a mobile app}

\mcqitem{6}{What does the formula $P = F \times S \times C$ calculate in user feedback prioritization?}{Product sales profit}{Priority score based on Frequency, Severity, and Confidence in evidence}{Software processing speed}{Project manager salary}

\mcqitem{7}{In user feedback prioritization, why does a rare safety failure (High Severity, Low Frequency) often outrank a common cosmetic glitch (Low Severity, High Frequency)?}{Cosmetic issues cannot be fixed}{Severe safety or task-blocking failures carry catastrophic risk and must be remediated immediately}{Frequency is the only metric that matters in engineering}{Severity is ignored in modern product design}

\mcqitem{8}{What is "Regression Testing" during design refactoring and product iteration?}{Testing whether older users can operate the app}{Verifying that recent improvements or changes have not unintentionally broken previously working functionality}{Testing software on obsolete computers}{Reverting to the very first sketch}

\mcqitem{9}{What are the four stages of the continuous improvement PDCA cycle?}{Produce, Deliver, Calculate, Archive}{Plan, Do, Check, Act}{Pitch, Demonstrate, Close, Acquire}{Propose, Design, Construct, Authorize}

\mcqitem{10}{What is the primary rule of Audience-Centred Product Presentation?}{Use the exact same technical jargon regardless of who is listening}{Tailor the language, narrative focus, and evidence to the specific priorities of the audience (investors, users, or engineers)}{Speak as fast as possible to fit all information into the slides}{Avoid showing the prototype to the audience}

\mcqitem{11}{When presenting a product to executive business investors, the presentation should focus most on:}{Transistor gate schematics and code line counts}{Market viability, unit economics, customer acquisition, scalability, and ROI}{CSS stylesheet class names}{The personal biographies of all junior designers}

\mcqitem{12}{What does the acronym NABC stand for in innovation pitching?}{Network, Action, Budget, Capital}{Need, Approach, Benefits per costs, Competition / alternatives}{Novelty, Architecture, Branding, Customer}{Narrative, Audience, Benchmark, Conclusion}

\mcqitem{13}{In the NABC framework, what does the "Need" section address?}{The team's need for funding}{A specific, validated user pain point backed by real-world evidence}{The server hardware requirements}{The company's legal registration}

\mcqitem{14}{In the NABC framework, why is it critical to explicitly analyze "Competition and Alternatives"?}{To mock competitors publicly}{To demonstrate how your solution provides superior relative advantage compared to existing workarounds, competitors, and non-consumption}{Because investors require buying competitor shares}{To copy competitors' marketing slogans}

\mcqitem{15}{What is an "Elevator Pitch"?}{A presentation delivered inside a moving elevator}{A concise 30–60 second verbal summary that clearly explains the problem, solution, unique value, and a requested call-to-action}{A 2-hour keynote address}{A written technical dissertation}

\mcqitem{16}{What is the ideal duration of a standard elevator pitch?}{10 to 15 seconds}{30 to 60 seconds}{15 to 20 minutes}{Exactly 1 hour}

\mcqitem{17}{Which of the following represents the most compelling opening "Hook" for a product presentation?}{"Hello, our names are on slide 1, and today we will talk about our project."}{"Every year, 40\% of critical hospital medicine spoils due to undetected temperature spikes in transit. Here is how we solve that."}{"We spent 100 hours writing 5,000 lines of Python code."}{"Our company was founded in 2026."}

\mcqitem{18}{What is the "Rule of One" in high-impact presentation slide design?}{Every presentation must only have one single slide}{Each slide should focus on conveying ONE core idea with clean visual support}{Only one person in the audience is allowed to ask questions}{Use only one color across all slides}

\mcqitem{19}{Why should product demonstrations focus on a "Representative User Task" rather than showing every single setting?}{Setting menus are confidential}{Demonstrating an authentic end-to-end task proves how the product delivers core value in real context}{It prevents the computer from crashing}{Users never use settings menus}

\mcqitem{20}{What is a "Design Rationale"?}{An excuse made when a project fails}{The explicit justification and evidence-based reasoning behind why a specific design choice was made}{The total monetary cost of the design}{The list of software libraries used}

\mcqitem{21}{Why is stating known limitations and failure conditions important during a final technical case study presentation?}{It makes the team look incompetent}{It demonstrates technical integrity, thorough testing, and builds strong credibility with evaluators}{Limitations are required by criminal law}{It guarantees immediate venture capital funding}

\mcqitem{22}{In a design case study, what does "Requirement Traceability" demonstrate?}{That code was written by tracing paper blueprints}{A clear, unbroken chain connecting user needs $\rightarrow$ engineering specifications $\rightarrow$ prototype features $\rightarrow$ verification test results}{Tracking the GPS coordinates of team members}{Calculating software download speeds}

\mcqitem{23}{When giving a live prototype demonstration, why is having a pre-recorded backup video considered essential?}{Because live prototypes are illegal in formal presentations}{To safeguard against unexpected live network, hardware, or environmental glitches during the pitch}{So the presenter can leave the room early}{Because video editing is mandatory for grading}

\mcqitem{24}{Which elevator pitch template is widely considered standard in product management?}{"For [target user] who [problem], our [product] is a [category] that [benefit]. Unlike [competition], we [differentiator]."}{"Buy our product because it is the cheapest and best in the world."}{"We have built a mobile app that does everything for everyone."}{"Our product has AI and blockchain built inside."}

\mcqitem{25}{What is the danger of relying on "Feature Bloat" (adding endless features) to save a struggling product?}{It increases software file size only}{It adds cognitive clutter, increases maintenance costs, and fails to fix an unvalidated core value proposition}{It makes the app too easy to use}{Feature bloat is always recommended by designers}

\mcqitem{26}{What does the "Benefits per Costs" section in the NABC model evaluate?}{The retail price of the prototype box}{The quantitative and qualitative value delivered compared to the financial, time, and adoption effort required by the user}{The monthly salary of the engineers}{The cost of office furniture}

\mcqitem{27}{During a final team product presentation, what is the best practice for team delivery?}{Have one person talk for 100\% of the time while others look at their phones}{Coordinate defined speaker transitions, maintain consistent terminology, and show shared ownership of the project}{Argue with teammates on stage about technical details}{Read directly from full paragraphs of text on the slides}

\mcqitem{28}{What should immediately follow the conclusion of an effective product pitch?}{A clear, specific "Ask" or Call-to-Action (e.g., requesting pilot testing access, partnership, or seed funding)}{Leaving the stage in total silence}{Apologizing for any mistakes made during the presentation}{Handing out paper business cards without speaking}

\mcqitem{29}{What type of strategic pivot occurred when YouTube pivoted from a video dating platform to a universal video hosting platform?}{Technology architecture pivot}{Problem / Value proposition \& customer segment pivot}{Zoom-in pivot}{Channel pivot}

\mcqitem{30}{What constitutes the final comprehensive deliverable of an INT335 Design Thinking innovation project?}{A single PowerPoint slide with a logo}{An evidence-backed portfolio including research synthesis, iterative prototypes, test logs, technical specs, and a persuasive demonstration}{A 100-page theoretical essay with no prototype}{A list of marketing slogans}


\newpage
\section*{Answer Key \& Step-by-Step Explanations}
\addcontentsline{toc}{section}{Answer Key \& Explanations}

\begin{answerkeytable}{Unit 6: Quick Answer Key}
\begin{tabular}{c c c c c c c c c c}
\toprule
\textbf{Q1} & \textbf{Q2} & \textbf{Q3} & \textbf{Q4} & \textbf{Q5} & \textbf{Q6} & \textbf{Q7} & \textbf{Q8} & \textbf{Q9} & \textbf{Q10} \\
\textbf{B} & \textbf{B} & \textbf{B} & \textbf{B} & \textbf{B} & \textbf{B} & \textbf{B} & \textbf{B} & \textbf{B} & \textbf{B} \\
\midrule
\textbf{Q11} & \textbf{Q12} & \textbf{Q13} & \textbf{Q14} & \textbf{Q15} & \textbf{Q16} & \textbf{Q17} & \textbf{Q18} & \textbf{Q19} & \textbf{Q20} \\
\textbf{B} & \textbf{B} & \textbf{B} & \textbf{B} & \textbf{B} & \textbf{B} & \textbf{B} & \textbf{B} & \textbf{B} & \textbf{B} \\
\midrule
\textbf{Q21} & \textbf{Q22} & \textbf{Q23} & \textbf{Q24} & \textbf{Q25} & \textbf{Q26} & \textbf{Q27} & \textbf{Q28} & \textbf{Q29} & \textbf{Q30} \\
\textbf{B} & \textbf{B} & \textbf{B} & \textbf{A} & \textbf{B} & \textbf{B} & \textbf{B} & \textbf{A} & \textbf{B} & \textbf{B} \\
\bottomrule
\end{tabular}
\end{answerkeytable}

\begin{expbox}[Detailed Pedagogical Explanations]
\begin{enumerate}[leftmargin=6mm, itemsep=2.5mm]
  \item \textbf{[Q1: Option B]} Iteration is the continuous learning loop of building, testing, learning, and refining.
  \item \textbf{[Q2: Option B]} Refactoring improves internal quality/structure; re-design overhauls the fundamental solution.
  \item \textbf{[Q3: Option B]} Invalidation of foundational assumptions or critical safety/architectural flaws mandates re-design.
  \item \textbf{[Q4: Option B]} A pivot is a structured hypothesis pivot that tests a new strategy while preserving past learning.
  \item \textbf{[Q5: Option B]} Moving from B2C students to B2B enterprise clients is a customer segment pivot.
  \item \textbf{[Q6: Option B]} $P = F \times S \times C$ ranks feedback by Frequency, Severity, and Confidence.
  \item \textbf{[Q7: Option B]} High-severity failures can cause total task abandonment or safety disasters.
  \item \textbf{[Q8: Option B]} Regression testing ensures new code or design changes do not break existing working flows.
  \item \textbf{[Q9: Option B]} The continuous improvement cycle is Plan $\rightarrow$ Do $\rightarrow$ Check $\rightarrow$ Act.
  \item \textbf{[Q10: Option B]} Effective communication aligns message and evidence directly with audience interests.
  \item \textbf{[Q11: Option B]} Investors focus on market viability, business sustainability, unit economics, and return.
  \item \textbf{[Q12: Option B]} NABC stands for Need, Approach, Benefits per costs, and Competition/alternatives.
  \item \textbf{[Q13: Option B]} "Need" defines the genuine, validated user pain point backed by research evidence.
  \item \textbf{[Q14: Option B]} Comparing with alternatives highlights the unique relative advantage of your innovation.
  \item \textbf{[Q15: Option B]} An elevator pitch delivers a high-impact product summary within 30–60 seconds.
  \item \textbf{[Q16: Option B]} Standard elevator pitches are strictly timeboxed between 30 and 60 seconds.
  \item \textbf{[Q17: Option B]} A compelling hook uses a vivid, quantified problem statement to command immediate attention.
  \item \textbf{[Q18: Option B]} The "Rule of One" limits each slide to one dominant idea to maximize audience retention.
  \item \textbf{[Q19: Option B]} Demonstrating a realistic user task proves the solution's real-world value proposition.
  \item \textbf{[Q20: Option B]} Design rationale provides the empirical evidence justifying why specific choices were made.
  \item \textbf{[Q21: Option B]} Acknowledging limitations demonstrates rigor, maturity, and builds evaluator trust.
  \item \textbf{[Q22: Option B]} Traceability connects user needs through engineering specs to test verification results.
  \item \textbf{[Q23: Option B]} A backup demo video protects against live hardware, venue Wi-Fi, or software crashes.
  \item \textbf{[Q24: Option A]} The standard positioning template clearly specifies user, problem, category, benefit, and differentiator.
  \item \textbf{[Q25: Option B]} Feature bloat increases interface complexity without addressing core usability or value problems.
  \item \textbf{[Q26: Option B]} Benefits per costs assesses total net value relative to the customer's effort, price, and friction.
  \item \textbf{[Q27: Option B]} Coordinated transitions and shared delivery reflect strong professional team cohesion.
  \item \textbf{[Q28: Option A]} Every presentation must culminate in a concrete, actionable call-to-action.
  \item \textbf{[Q29: Option B]} YouTube pivoted its core value proposition and audience from video dating to universal video sharing.
  \item \textbf{[Q30: Option B]} A complete design project integrates research, prototypes, test logs, technical specs, and pitch.
\end{enumerate}
\end{expbox}
